\chapter{Simulación y validación de la cadena de procesamiento}

En este capítulo se presentan resultados de simulación obtenidos con la cadena de procesamiento descrita en el capítulo anterior. El objetivo es validar el comportamiento de cada bloque —generación de escenario, filtrado de clutter, cálculo de la función de ambigüedad cruzada y detección CFAR— en un entorno controlado, verificando que los resultados obtenidos coincidan con los valores físicos esperados.

El flujo general de simulación se muestra en la Figura~\ref{fig:flow_pipeline_sim}. En primer lugar se genera un escenario bistático con un transmisor, un receptor, un blanco móvil y múltiples reflectores estacionarios que representan el clutter. A partir de esta geometría se construyen las señales de referencia y vigilancia, las cuales son procesadas por la cadena digital.

\begin{figure}[H]
	\centering
	\includegraphics[width=0.85\textwidth]{pr/flow_pipeline_sim.pdf}
	\caption{Diagrama de flujo del proceso de simulación y procesamiento.}
	\label{fig:flow_pipeline_sim}
\end{figure}

\section{Escenario de simulación}

El escenario considerado consiste en un sistema de radar pasivo bistático con un transmisor fijo, un receptor fijo y un blanco en movimiento. Además, se incorporan múltiples reflectores estacionarios distribuidos en el entorno para modelar el clutter.

La Figura~\ref{fig:geometry_sim} muestra la geometría del escenario simulado.

\begin{figure}[H]
	\centering
	\includegraphics[width=0.6\textwidth]{sim_data/fig01_geometry.png}
	\caption{Geometría del escenario de simulación.}
	\label{fig:geometry_sim}
\end{figure}

A partir de esta configuración geométrica se generan las señales de referencia y vigilancia. La señal de vigilancia incluye la contribución del eco del blanco, la trayectoria directa del transmisor, múltiples reflexiones estacionarias y ruido aditivo.

El simulador devuelve, además, los valores de referencia del blanco (``ground truth''), en particular el rango bistático y el corrimiento Doppler esperados. Estos valores se utilizan posteriormente para validar la detección obtenida mediante la CAF y el detector CFAR.

\section{Análisis espectral de las señales}

Antes de aplicar la cadena de procesamiento, resulta útil analizar el contenido espectral de las señales generadas. Para ello se calcula la densidad espectral de potencia (PSD) utilizando el estimador implementado en el método \texttt{plot\_psd} del pipeline.

La Figura~\ref{fig:psd_input} muestra la PSD de las señales de entrada al sistema, es decir, el canal de referencia y el canal de vigilancia.

\begin{figure}[H]
	\centering
	\includegraphics[width=0.8\textwidth]{sim_data/fig02_psd_input.png}
	\caption{PSD de las señales de entrada al pipeline.}
	\label{fig:psd_input}
\end{figure}

Puede observarse que ambas señales comparten el mismo ancho de banda, ya que la señal de vigilancia contiene copias retardadas y escaladas de la señal de referencia. Sin embargo, el canal de vigilancia presenta mayor potencia debido a la presencia de la trayectoria directa y del clutter, componentes altamente correlacionados con la referencia.

Este comportamiento es consistente con el modelo de señal presentado en el capítulo anterior, donde la interferencia de trayectoria directa y el clutter dominan la energía de la señal recibida.

\section{Efecto del filtro de clutter}

El siguiente paso consiste en aplicar el filtro adaptativo \textit{block lattice} para cancelar la trayectoria directa y las reflexiones estacionarias.

La Figura~\ref{fig:psd_caf_in} muestra la PSD de las señales que ingresan al bloque de cálculo de la CAF, es decir, la referencia y la señal de vigilancia luego del filtrado de clutter.

\begin{figure}[H]
	\centering
	\includegraphics[width=0.8\textwidth]{sim_data/fig03_psd_caf_in.png}
	\caption{PSD a la entrada del bloque CAF luego del filtrado de clutter.}
	\label{fig:psd_caf_in}
\end{figure}

Se observa una reducción significativa de la potencia asociada a las componentes correlacionadas con la referencia, mientras que el ancho de banda de la señal se mantiene inalterado. Este comportamiento es el esperado para un filtro adaptativo de cancelación de clutter, cuyo objetivo es eliminar únicamente las componentes correlacionadas con la señal de referencia sin distorsionar la información restante.

Desde un punto de vista geométrico, este resultado refleja el proceso de proyección ortogonal descrito en la sección del filtro lattice: las componentes asociadas al subespacio generado por las copias retardadas de la referencia son removidas, preservando las componentes no correlacionadas, entre ellas los ecos de blancos en movimiento.

\section{Función de ambigüedad cruzada}

Una vez filtrada la señal de vigilancia, se calcula la función de ambigüedad cruzada mediante el algoritmo \textit{batch}. Este proceso produce un mapa rango--Doppler en el cual los blancos se manifiestan como máximos locales.

La Figura~\ref{fig:caf_no_cf} muestra la CAF obtenida sin aplicar el filtro de clutter.

\begin{figure}[H]
	\centering
	\includegraphics[width=0.8\textwidth]{sim_data/fig04_caf_no_clutter_filter.png}
	\caption{CAF sin filtrado de clutter.}
	\label{fig:caf_no_cf}
\end{figure}

En este caso, la presencia de la trayectoria directa y del clutter genera un nivel elevado de energía alrededor del origen Doppler, lo cual dificulta la identificación del blanco.

La Figura~\ref{fig:caf_cf} muestra la CAF obtenida luego de aplicar el filtro de clutter.

\begin{figure}[H]
	\centering
	\includegraphics[width=0.8\textwidth]{sim_data/fig05_caf_with_clutter_filter.png}
	\caption{CAF luego del filtrado de clutter.}
	\label{fig:caf_cf}
\end{figure}

Puede observarse una mejora significativa en el contraste del mapa rango--Doppler, lo cual permite identificar con mayor claridad el eco del blanco. Este resultado confirma la utilidad del filtrado adaptativo para aumentar el rango dinámico efectivo del sistema.

\section{Detección mediante CFAR}

La detección de blancos se realiza aplicando el detector CA-CFAR sobre la magnitud de la CAF. Para analizar el comportamiento del detector, se considera un corte Doppler en el bin de rango correspondiente al blanco simulado.

La Figura~\ref{fig:cut_no_cf} muestra este corte sin filtrado de clutter, mientras que la Figura~\ref{fig:cut_cf} muestra el resultado luego del filtrado.

\begin{figure}[H]
	\centering
	\includegraphics[width=0.8\textwidth]{sim_data/fig06_cut_no_clutter_filter.png}
	\caption{Corte Doppler sin filtrado de clutter.}
	\label{fig:cut_no_cf}
\end{figure}

\begin{figure}[H]
	\centering
	\includegraphics[width=0.8\textwidth]{sim_data/fig07_cut_with_clutter_filter.png}
	\caption{Corte Doppler luego del filtrado de clutter.}
	\label{fig:cut_cf}
\end{figure}

En el caso filtrado, el pico correspondiente al blanco supera claramente el umbral de detección estimado por el algoritmo CFAR. Este comportamiento es consistente con la teoría presentada en el capítulo anterior: al reducir la potencia del clutter, la estimación local del ruido disminuye y el detector puede operar con mayor sensibilidad.

\section{Validación}

La posición del máximo detectado en la CAF coincide con los valores de rango bistático y frecuencia Doppler calculados por el simulador. Esto confirma la correcta implementación de la cadena de procesamiento completa.

En conjunto, los resultados obtenidos validan el funcionamiento del sistema de radar pasivo simulado y muestran la importancia del filtrado de clutter para mejorar la detectabilidad de blancos en escenarios con interferencias dominantes.
